%%%%%%%%%%%%%%%%%%%%%%%%%%%%%%%%%%%%%%%%%%%%%%%%%%%%%%%%%%%%%%%%%%%%%%
% How to use writeLaTeX:
%
% You edit the source code here on the left, and the preview on the
% right shows you the result within a few seconds.
%
% Bookmark this page and share the URL with your co-authors. They can
% edit at the same time!
%
% You can upload figures, bibliographies, custom classes and
% styles using the files menu.
%
% If you're new to LaTeX, the wikibook is a great place to start:
% http://en.wikibooks.org/wiki/LaTeX
%
%%%%%%%%%%%%%%%%%%%%%%%%%%%%%%%%%%%%%%%%%%%%%%%%%%%%%%%%%%%%%%%%%%%%%%
\documentclass{tufte-handout}

%\geometry{showframe}% for debugging purposes -- displays the margins

\usepackage{amsmath}

% Set up the images/graphics package
\usepackage{graphicx}
\setkeys{Gin}{width=\linewidth,totalheight=\textheight,keepaspectratio}
\graphicspath{{graphics/}}

\title{Stata Lab 2: Data Management}
\author{DIME Analytics \\ dimeanalytics@worldbank.org}
% \date{24 January 2009}  % if the \date{} command is left out, the current date will be used

% The following package makes prettier tables.  We're all about the bling!
\usepackage{booktabs}

% The units package provides nice, non-stacked fractions and better spacing
% for units.
\usepackage{units}

% The fancyvrb package lets us customize the formatting of verbatim
% environments.  We use a slightly smaller font.
\usepackage{upquote}
\usepackage{fancyvrb}
\fvset{fontsize=\normalsize}
\renewcommand{\FancyVerbFormatLine}{\color{violet}}
\DefineShortVerb{\|}

% Small sections of multiple columns
\usepackage{multicol}

% Provides paragraphs of dummy text
\usepackage{lipsum}

% These commands are used to pretty-print LaTeX commands
\newcommand{\doccmd}[1]{\texttt{\textbackslash#1}}% command name -- adds backslash automatically
\newcommand{\docopt}[1]{\ensuremath{\langle}\textrm{\textit{#1}}\ensuremath{\rangle}}% optional command argument
\newcommand{\docarg}[1]{\textrm{\textit{#1}}}% (required) command argument
\newenvironment{docspec}{\begin{quote}\noindent}{\end{quote}}% command specification environment
\newcommand{\docenv}[1]{\textsf{#1}}% environment name
\newcommand{\docpkg}[1]{\texttt{#1}}% package name
\newcommand{\doccls}[1]{\texttt{#1}}% document class name
\newcommand{\docclsopt}[1]{\texttt{#1}}% document class option name

\begin{document}

\maketitle% this prints the handout title, author, and date

\begin{marginfigure}%
  \includegraphics[width=\linewidth]{light.png}
\end{marginfigure}

\begin{abstract}
In this lab we will use our data management skills to complete two common, but complex tasks: reshaping data and merging master data.

This exercise is challenging and will probably only be finished by those of you already familiar with some of the topics covered. This exercise is written in a step-by-step recipe fashion, so it should be possible to finish the task on your own if you do not finish it today.

For those of you who are very familiar with these tasks and complete them quickly, there is a bonus exercise that introduces a helpful new command for data cleaning, |iecodebook|.

\bigskip\noindent \textbf{Exercise Objectives}:
\begin{enumerate}
  \item Get used to the folder structure and load the raw data
  \item Introduction to creating master data sets
  \item Introduction to how to address unexpected values
  \item Load “wide” data and |reshape| it into “long” data
  \item Use a master dataset to |merge| data
\end{enumerate}
\end{abstract}

%\printclassoptions
\section{Part 1: Get used to the folder structure (again)}

As you know from Lab 2, the |iefolder| command sets up
the basic folder structure for all data work for an entire project.
But it is up to its users to make use of
and maintain the basic structure so that
new users can intuitively understand your work.
It is often the case that research assistants
are the ones leading in making sure this is accomplished;
however, it is important that field coordinators are familiar with this system
so you can understand what RAs have done –
since you will often work with two, three, or more RAs
over the course of a single project!

As you work with full projects,
you will inevitably get more used to the structure,
but here is an introduction exercise to get you started.
If you have not yet created a |/Lab3/| folder,
do so now using |iefolder|.
In the first lab, we did this using:
\begin{Verbatim}
  iefolder new subfolder Lab1, project("${directory}")
\end{Verbatim}
where |directory| was the name of the global macro
we created to store the location of our project folder.

If you have closed Stata since then,
you will need to run code to creates that macro again.
All global macros are cleared at the end of the session,
and are also cleared whenever we begin a session with |ieboilstart|.
That is also why this segment of code appears in |Project_MasterDofile.do|:
\begin{Verbatim}[frame=leftline,numbers=left]
*Set this value to the user currently using this file
  global user  1

* Root folder globals
* ---------------------

  if $user == 1 {
    global projectfolder "/Users/bbdaniels/GitHub/DIME-MSIE-Workshop/Material/Labs/Stata-Track-2"
  }
  if $user == 2 {
    global projectfolder ""  // Enter the file path to the project folder for the next user here
  }
\end{Verbatim}
This section of code allows you to change between
users of the code, by pre-specifying where the project folder
is located on each user's computer.
Since the project folder will be identical for all users,
the remainder of the code creates additional global macros
that hold that information separately:
\begin{Verbatim}[frame=leftline,numbers=left]
 *Lab1 sub-folder globals
 global Lab1                   "${dataWorkFolder}/Lab1"
 global Lab1_encrypt           "${encryptFolder}/Subfolder Lab1 Encrypted"

 *Lab2 sub-folder globals
 global Lab2                   "${dataWorkFolder}/Lab2"
 global Lab2_encrypt           "${encryptFolder}/Subfolder Lab2 Encrypted"

 *Lab3 sub-folder globals
 global Lab3                   "${dataWorkFolder}/Lab3"
 global Lab3_encrypt           "${encryptFolder}/Subfolder Lab3 Encrypted"
\end{Verbatim}
...and so on. It is now straightforward for multiple users
to share the entire folder structure,
because allowing another person to work
on the same code is as simple as adding their filepath
to the ``root folder globals'' section.
Notice that, because all the global macros are fully expanded,
if you write
\begin{Verbatim}
  di "${Lab2}"
\end{Verbatim}
Stata will show you the complete filepath to the folder.

As you work with full projects,
instead of using subfolders,
you will use |round| folders instead.
You can try creating one as a demo using a command like:
\begin{Verbatim}
  iefolder new round Baseline, project("${projectfolder}")
\end{Verbatim}
Note that we can now use the |projectfolder| global in the code.
Take a minute to explore the Baseline folder,
and see what global definitions were added tp |Project_MasterDofile.do|.
In a full project, this file stucture is used to stay organized,
with data and code in the correct place for various tasks.
(In these exercises, we will continue using the simpler subfolders
to get used to using a single master file and globals.)

\section{Part 2: Create a master dataset}

If you have not done so yet, create a Lab 2 subfolder using:
\begin{Verbatim}
  iefolder new subfolder Lab2, project("${projectfolder}")
\end{Verbatim}
Copy two datasets into that folder:
|endline_data_raw_nodup.dta|, which we provided, and
|endline_data_final.dta|, which we created in Lab 1.
(There is a spare copy for you in the Materials folder.)
Now we are ready to get to work.

First, we are going to create a household master data set
from the data we created in Lab 1.
Start a new do-file and save it as:
|/Lab3/Dofiles/Construct/construct-masterhh.do.|
Use the dataset |endline_data_final.dta| that we created.
Do |rename id_05 id| and check that it is an ID variable using |isid|.

Label the household head variables |pl_*_1| so they will be interpretable,
then use |rename| and |label| to mark everything as household-level variables
(Hint: try |rename * hh_*|). Then |order| the id variables first,
putting everything else in sequential order (not alphabetical),
|compress| the data, and |saveold| this dataset as
|master_households.dta| in the |/MasterData/| folder
(using the global macro created in |Project_MasterDofile.do|)
The |order| command makes the dataset easy to understand
for the next person who opens it
and |compress| ensures that it takes as little space as possible
on the computer (and in transit).
This might look like the one below.

\section{Part 3: Clean, reshape, and merge a roster}

\subsection{Identify unexpected values in data}

In most cases we have research assistants cleaning the data set.
But there are parts of the cleaning process
that field coordinators usually are at least involved in resolving.
This section will focus on those tasks.
The most important task related to unexpected values
for both research assistants and field coordinators is to identify them.
In the end, it is up to the researcher in charge of the project
to make the final decision on how to address them.
The main lesson in this exercise is therefore to learn
methods for exploring the data set for unexpected data values.
But we will cover some examples on how to address them as well.

Open Stata and a blank dofile (in the same dofile editor).
Now, |use| the dataset |endline_data_raw_nodup.dta|, in the do-file.
Let’s save this do-file now call it |reshape-roster.do|.
As you know, Stata does not auto-save,
so make sure to save your do-file frequently.
(Again, it is even easier if you do this in a separate editor.)
As with any data set you open for the first time,
you should explore the data set to make sure that you understand
the unit of observation and have a uniquely and fully identifying ID variable.
For this exercise you do not have to correct anything you find,
since we don't have field notes to go back to,
but it is best to go through the process of exploring the data set
to get into this extremely important habit.

We are going to start by trimming this dataset
down to the household roster and working from there.
First, |keep| only the household ID variable and the
|age_?|, |pl_sex_?|, and |education_?| variables from the household roster.
Then, use the command |codebook, compact| on all variables.
If you do not specify any variable Stata will assume that you mean all variables.
What do you notice immediately?
Look first at the columns |Obs| and |Unique|.
For the ID variable, these should both be equal,
and equal to the number of observations in the data set.
Here you can see they are. If they were not,
you could use |duplicates report id_05|
to see how many errors there are and |duplicates list id_05| to see them.
The |duplicates| command has many more features and is also worth exploring!

For other variables, we are more concerned about missing values.
Where the number in the |Obs| column match the number of observations in the data,
we have no missing values.
If the number is lower than the number of observations,
we have observations with missing data.
Sometimes that is ok, for example, if a question was not asked
because of a skip pattern.
Since you are not familiar with the questionnaire for this particular dataset,
we will not ask you to follow up on missing values in this exercise,
but you should do that in a real-life scenario.
This data set is not as large as most survey datasets,
so you can scroll through all variables.
In a real data set, that would be difficult and you could, for example,
limit the summarize command to the education variables by typing
|codebook education_*, compact|.
Look at the max and min values for the education variables.
Are any of these values unexpected? Are any of the values too low or too high?

Negative values do not make sense for an education variable,
so -88 must be a survey code for a type of missing value.
For example, “don’t know” or “declined to answer” could be coded like this.
Type |codebook education_1|.
In general, you will consult the questionnaire or
the person who wrote the questionnaire to find out what the code represents.
The |summarize| command also has an option called |detail|.
Enter |summarize education_1, detail| in Stata.
The first column with the title |Percentiles|
gives you the value of some standard percentiles.
Most observations in this data set did not have any income from renting out land.
The middle column shows the four smallest values and the four largest values.
The third column shows additional descriptive statistics.
Another method to identify unexpected values is |histogram|.
Histograms are particularly useful to explore outliers.
Type |histogram education_1| in Stata,
and you can see immediately where the data is unusual.

\subsection{Reshaping the household roster dataset}

The variables are right now in what we call a “wide” format:
we have a dataset that (in our minds) corresponds to individuals,
but each observation is actually a household.
Therefore the individuals are listed across the rows (“wide”),
and we would like to |reshape| the dataset so each row represents an individual.
Type |help reshape| and give it a shot. (It is not going to work.)

We need to do a little cleaning so that the |reshape| can proceed.
Let’s rename some variables so they are more sensible
using |rename|’s powerful syntax.
Let’s try to have our dataset en up only with
|age?|, |sex?|, and |edu?| variables using just three lines of code.
(Hint: see what |rename *sex_* sex*[2]| does.
Then |rename id_05 hh_id|.
This is so that it matches exactly the corresponding code
in the master dataset you just created.

Now we have a dataset |reshape| can work with. Use:
\begin{Verbatim}
  reshape long age sex edu, i(hh_id) j(mem_id)
\end{Verbatim}
Then |browse| the data. How’s it look?
Ok, no, it’s not great. Raw data never is.
Clean and label this data fully –
including, at this point, using |drop if (sex == .)|.
Dropping observations with key variables missing
does clean up after |reshape|,
but in this case you’ll notice that some households
are going to be dropped entirely: they have no members with data.
Normally we’d go back and ask the field team why this is the case,
but this will be used as an example for now.

In part 2, we discussed how to identify codes
used to represent missing values and outliers.
You can also do that now that |reshape| has condensed the variables.
To see what we are working with, do the following:
type |codebook edu| and see that the name of
the associated value label is |ed_level|.
Then type |numlabel ed_level, add| and finally,
|tab edu, missing plot|. Nearly magical, right?

Each code -77, -88, ..., means something and
if we would change them all to the same missing value (.),
we would lose information.
The researcher might want to know how many
of the respondents that did not answer the question did it
because they “don’t know”, “declined to answer” or the “question do not apply”.
Stata provides a solution in its extended missing values (.a, .b, .c ... .z).
If you are not familiar with missing values you can type |help missing|
in Stata and read the documentation.
We now need to replace the survey codes with these missing values.
Do this now for |edu|, and add labels to |.a| and |.b|.
Finally, |save| the data.
When you are done, your dofile might look like the one below.
There is a trick you can use to easily loop over all numeric variables
if there are more than one using this scheme,
but you will not need to do this here.
Use |ds, has(type numeric)| to generate a varlist of all numeric variables
that you then can loop through.
If you do not know how to use |ds|, type |help ds| in Stata.

\subsection{Merge the master data and the roster data}

Start a new dofile and save it as |construct-roster.do|.
This will be a short dofile: all we are going to do is |merge| on our master data.
First, |use| the clean roster data you saved.
Then type |help merge| and think about how to merge
|master_households.dta| onto the roster dataset we’ve created. Do that.
Once you do, |tab| the special |_merge| variable that is created.
What does it tell you? Next, |drop| the |_merge| variable.
Future merges will fail if the dataset contains this variable.
Then, create a “tag” variable called |tag_hh|
indicating the first person in each household using |egen| and label it.
This will allow you to continue conducting analysis at the household level
even in the combined dataset.
Use |isid| to confirm that you have a unique ID variable (or combination).
If you don’t, can you find out what the problem is?
(Just for now, go ahead and |drop| the offending observations.)
Finally, save the dataset as |hh_roster.dta|.
When you are done, your dofile might look like the one below.

\subsection{Using Excel to generate repetitive code}

You can use Excel to manage repetitive code.
A lot of what we have done in terms of cleaning
could be done faster this way using a program called |iecodebook|,
available at \url{http://worldbank.github.io/iefieldkit}.
See if you can do some cleaning and labeling quickly using this tool.

\begin{figure*}[h]
{\setstretch{0.7}
\VerbatimInput[frame=lines,numbers=left,label=construct-masterhh.do]
{../DataWork/Lab2/construct-masterhh.do}}
\end{figure*}


\begin{figure*}[h]
{\setstretch{0.7}
\VerbatimInput[frame=lines,numbers=left,label=reshape-roster.do]
{../DataWork/Lab2/reshape-roster.do}}
\end{figure*}


\begin{figure*}[h]
{\setstretch{0.7}
\VerbatimInput[frame=lines,numbers=left,label=construct-roster.do]
{../DataWork/Lab2/construct-roster.do}}
\end{figure*}

Thanks to the template creators.\cite{tuftelatex}

\bibliography{sample-handout}
\bibliographystyle{plainnat}

\end{document}
